\subsection{Related Work on Global Robustness}

While \cite{diff} does not tackle global robustness per se, it is one of the first works to consider {\em differential} verification, introducing the {\em diff} variables that are used here and partly in \cite{ITNE}. 
The paper \cite{diff} did not consider MILP encoding.
%
Some works \cite{Leino,Zhang22a,Sun22,Chen21,REGLO} focus on training a network to be, rather than proving the DNN is, globally robust.
%A key aspect of global robustness involves ensuring that analyzed data is both semantically meaningful and in-distribution (ID). 
Our approach to obtain DNN with better global robustness bounds uses (linear) auto encoders~\cite{hinton2006reducing}.
The (L)AE-DNN architectures (Fig.~\ref{fig:AE-DNN} in Section \ref{sec.pca}) projects the input on a dimensionaly reduced space to filter non-ID features 
\cite{aggarwalOutlierAnalysis2017}.

Some works \cite{Bastani16,Ruan19,Gopinath18} estimate global robustness bounds, but without hard verification certificates, e.g., using probabilistic guarantees \cite{Levy23,Mangal19}.
%More related, s
Several works \cite{Marabou,lipshitz,ITNE,GROCET} consider certifying {\em restricted} Lipschitz bound, to understand how the output can change given an input perturbation. However, they do not consider whether the classification changes. 
%
Closest to our work is VHAGaR \cite{vhagar}, which computes 
bounds $\bar{\alpha}^\varepsilon_{i,j}$ similar to our bounds 
$\bar{\beta}^\varepsilon_{i,j}$.
%, but limited to $L_\infty$-perturbations and variants. 
Crucially, VHAGaR bounds are not small enough for real-time robustness verification, whereas our bounds are. 
%One reason is {\em Diff} MILP is more accurate in practice.

Technically, none of the MILP-based works \cite{vhagar,lipshitz,ITNE} consider $L_1$-perturbations. Further, they use the classical MILP encoding from \cite{MILP}.
%, and not our novel {\em Diff} model, nor even the "1b" model. 
Note that \cite{ITNE} considers the {\em diff variables} $y,\hat{y}$, adding explicitly one {\em linear} constraint (Eq. (3) in \cite{ITNE})
%$\frac{y_i-\gamma_i}{2} \leq \hat{y}_i \leq \frac{y_i+\gamma_i}{2}$ 
implied by each of our three MILP models {\em Diff}, "2b+1" and "1b" models. \cite{vhagar} does things differently, adding per neuron constraints (depending on the perturbation) between the (binary) variables to simplify the MILP problem, which is well-adapted to occlusion and patches but less to $L_1$-perturbations. 
%(and to some extent to $L_\infty$-perturbations). 


%While non-linear AEs can achieve better compression and noise reduction, we consider linear AE in this work to avoid the additional computational complexity associated with non-linear activation functions within the MILP framework~\cite{fischetti2018deep}.

